Induct on the number of generators $x_1,\ldots,x_n$. If they are algebraically independent, there is nothing to prove. Otherwise choose a nonzero polynomial $f$ with $f(x_1,\ldots,x_n)=0$, and let $F$ be its highest homogeneous part.

Suppose first that $k$ is infinite. Choose $\lambda_1,\ldots,\lambda_{n-1}$ with $F(\lambda_1,\ldots,\lambda_{n-1},1)\neq0$, and put $x_i'=x_i-\lambda_i x_n$. After substitution, the coefficient of the highest power of $x_n$ is this nonzero scalar. Dividing by it gives a monic equation for $x_n$ over $A'=k[x_1',\ldots,x_{n-1}']$. Thus $A=A'[x_n]$ is integral over $A'$. Apply induction to $A'$ and transitivity of integrality. The resulting $y_i$ are linear combinations of the original generators.

For arbitrary $k$, choose positive integers $N_1,\ldots,N_{n-1}$ such that the weights $\sum_{i<n}N_i a_i+a_n$ of the monomials occurring in $f$ are distinct. For example, use successive powers of an integer larger than every exponent in $f$. Set $x_i'=x_i-x_n^{N_i}$. In the substituted polynomial, the monomial of largest weight gives a unique highest power of $x_n$ with nonzero scalar coefficient. Again $x_n$ is integral over $k[x_1',\ldots,x_{n-1}']$, and induction proves the assertion.

For the geometric statement, lift the linearly chosen $y_i$ to linear forms on $k^n$. They are linearly independent, since they are algebraically independent in the coordinate ring, so they define a surjective linear map $k^n\to k^r$. Evaluation at any point of $k^r$ gives $k[y_1,\ldots,y_r]\to k$, which extends to $A\to k$ by \exref{5.2}. This extension is evaluation at a point of $X$. Hence $X\to k^r$ is surjective. Identifying $k^r$ with the image of a linear section gives the required linear subspace.
